\documentclass[11pt]{article}

\usepackage[letterpaper,margin=0.95in]{geometry}
\usepackage[T1]{fontenc}
\usepackage[charter]{mathdesign}
\usepackage[scaled=0.92]{helvet}
\usepackage{microtype}
\usepackage{amsmath}
\usepackage{bm}
\usepackage{booktabs,tabularx,array}
\usepackage{enumitem}
\usepackage{xcolor}
\usepackage{titlesec}
\usepackage{caption}
\usepackage{float}
\usepackage{graphicx}
\usepackage[hidelinks]{hyperref}

\definecolor{headblue}{RGB}{24,54,96}
\definecolor{notebg}{RGB}{243,246,250}
\definecolor{fixbg}{RGB}{252,240,232}
\setlength{\parindent}{0pt}
\setlength{\parskip}{5pt}
\setlist{nosep,leftmargin=1.6em,topsep=3pt,itemsep=3pt}
\renewcommand{\thesection}{\Alph{section}}
\renewcommand{\thesubsection}{\thesection.\arabic{subsection}}
\titleformat{\section}{\sffamily\large\bfseries\color{headblue}}{\thesection}{0.6em}{}
\titlespacing*{\section}{0pt}{16pt}{4pt}
\titleformat{\subsection}{\sffamily\normalsize\bfseries}{\thesubsection}{0.6em}{}
\titlespacing*{\subsection}{0pt}{10pt}{2pt}
\numberwithin{equation}{section}
\renewcommand{\thetable}{\Roman{table}}
\captionsetup[table]{name=TABLE,labelsep=newline,font={footnotesize,sc},justification=centering,skip=3pt}

\newcommand{\dstar}{d^{*}}
\newcommand{\ez}{\hat{\bm e}_z}
\newcommand{\ah}{\hat{\bm a}}
\newcommand{\R}{\mathbb{R}}
\newcommand{\tr}{\operatorname{tr}}
\newcommand{\note}[1]{\par\smallskip\noindent\colorbox{notebg}{\parbox{\dimexpr\linewidth-2\fboxsep}{\small #1}}\par\smallskip}
\newcommand{\fix}[1]{\par\smallskip\noindent\colorbox{fixbg}{\parbox{\dimexpr\linewidth-2\fboxsep}{\small\textbf{Correction.} #1}}\par\smallskip}
\newcommand{\meta}[2]{\par\smallskip{\small\raggedright\textsf{Check:} \texttt{#1}\quad\textsf{Code:} \texttt{#2}\par}}

% numbers generated by the scripts (never typed by hand)
\IfFileExists{../results/calib/tables.tex}{% generated by scripts/calibrate_contact.py -- do not edit
\newcommand{\calibsetups}{%
Menagerie as shipped & 0-17 & 0.0-2.6 & 0.134 & dropped & 0.0 & -- & 0.00 & 0 \\
1-point pad, MuJoCo default contacts & 1-1 & 61.3-64.4 & 6.158 & held & 9.0 & 2.25 & 2.50 & 34 \\
1-point pad (simulator check) & 1-1 & 70.0-70.0 & 0.061 & held & 0.0 & 4.50 & 4.75 & 10 \\
grid pad, rigid (hard plastic) & 16-16 & 70.0-70.0 & 0.007 & held & 8.6 & 7.75 & 8.00 & 1 \\
grid pad, rubber (study) & 16-16 & 70.0-70.0 & 0.176 & held & 3.1 & 7.25 & 7.50 & 0 \\
}
\newcommand{\calibgeneral}{%
point (rigid) & 0.5 & 70 & 6 & 0.6 & 4.55--4.74 & 4.58 (-3.3\%) & 4.65 (-1.9\%) \\
point (rigid) & 0.8 & 70 & 6 & 0.6 & 7.30--7.60 & 7.35 (-3.4\%) & 7.44 (-2.1\%) \\
point (rigid) & 0.3 & 70 & 6 & 0.6 & 2.71--2.83 & 2.73 (-3.3\%) & 2.78 (-1.5\%) \\
point (rigid) & 0.5 & 50 & 6 & 0.6 & 3.24--3.37 & 3.26 (-3.3\%) & 3.32 (-1.7\%) \\
point (rigid) & 0.5 & 30 & 6 & 0.6 & 1.92--2.00 & 1.94 (-3.2\%) & 1.97 (-1.3\%) \\
point (rigid) & 0.5 & 70 & 4 & 0.6 & 3.04--3.16 & 3.06 (-3.2\%) & 3.11 (-1.6\%) \\
point (rigid) & 0.5 & 70 & 8 & 0.6 & 6.07--6.32 & 6.11 (-3.3\%) & 6.19 (-2.1\%) \\
point (rigid) & 0.5 & 70 & 6 & 0.8 & 3.40--3.55 & 3.43 (-3.4\%) & 3.48 (-1.7\%) \\
grid, rubber & 0.5 & 70 & -- & 0.6 & 7.23--7.53 & 7.19 (-4.5\%)$^\dagger$ & 7.44 (-1.1\%) \\
grid, rubber & 0.3 & 70 & -- & 0.6 & 4.29--4.47 & 4.30 (-3.8\%) & 4.43 (-0.8\%) \\
grid, rubber & 0.4 & 70 & -- & 0.6 & 5.76--6.00 & 5.75 (-4.2\%)$^\dagger$ & 5.95 (-0.8\%) \\
grid, rubber & 0.5 & 40 & -- & 0.6 & 4.08--4.25 & 4.07 (-4.2\%) & 4.26 (+0.2\%)$^\dagger$ \\
grid, rubber & 0.5 & 55 & -- & 0.6 & 5.66--5.89 & 5.63 (-4.5\%)$^\dagger$ & 5.89 (-0.0\%) \\
grid, rubber & 0.5 & 70 & -- & 0.8 & 5.39--5.62 & 5.35 (-4.9\%)$^\dagger$ & 5.60 (-0.2\%) \\
grid, rubber & 0.6 & 70 & -- & 0.9 & 5.76--6.00 & 5.68 (-5.3\%)$^\dagger$ & 6.00 (-0.1\%) \\
grid, rigid & 0.5 & 70 & -- & 0.6 & 7.23--7.53 & 7.41 (-1.5\%) & 7.79 (+3.6\%)$^\dagger$ \\
grid, rigid & 0.3 & 70 & -- & 0.6 & 4.29--4.47 & 4.40 (-1.5\%) & 4.65 (+4.1\%)$^\dagger$ \\
grid, rigid & 0.4 & 70 & -- & 0.6 & 5.76--6.00 & 5.92 (-1.3\%) & 6.26 (+4.4\%)$^\dagger$ \\
grid, rigid & 0.5 & 40 & -- & 0.6 & 4.08--4.25 & 4.28 (+0.7\%)$^\dagger$ & 4.71 (+10.8\%)$^\dagger$ \\
grid, rigid & 0.5 & 55 & -- & 0.6 & 5.66--5.89 & 5.87 (-0.3\%) & 6.27 (+6.5\%)$^\dagger$ \\
grid, rigid & 0.5 & 70 & -- & 0.8 & 5.39--5.62 & 5.59 (-0.5\%) & 5.96 (+6.1\%)$^\dagger$ \\
grid, rigid & 0.6 & 70 & -- & 0.9 & 5.76--6.00 & 6.00 (-0.1\%) & 6.43 (+7.1\%)$^\dagger$ \\
}
\newcommand{\calibapeak}{0.401}
\newcommand{\calibpointpred}{4.74}
\newcommand{\calibgridpred}{7.53}
\newcommand{\calibgridpredpeak}{7.23}
\newcommand{\calibinbandonePoint}{8/8}
\newcommand{\calibinbandtenPoint}{8/8}
\newcommand{\calibinbandoneRubber}{2/7}
\newcommand{\calibinbandtenRubber}{6/7}
\newcommand{\calibinbandoneRigid}{6/7}
\newcommand{\calibinbandtenRigid}{0/7}
}{}
\IfFileExists{../results/calib/characterize_rubber.tex}{% generated by scripts/characterize_contact.py -- do not edit
\newcommand{\calibcharacterizerubber}{%
reference (box 20x4x4 cm) & 16 & 7.23--7.53 & 7.17 (-5\%)$^\dagger$ & 7.41 (-2\%) \\
box yawed 3 deg & 16 & 7.23--7.53 & 4.27 (-43\%)$^\dagger$ & 6.14 (-18\%)$^\dagger$ \\
box yawed 10 deg & 16 & 7.23--7.53 & 0.16 (-98\%)$^\dagger$ & 1.13 (-85\%)$^\dagger$ \\
grasp +1 cm toward the weight & 16 & 8.23--8.53 & 8.14 (-5\%)$^\dagger$ & 8.34 (-2\%) \\
grasp -1 cm away from the weight & 16 & 6.23--6.53 & 6.17 (-5\%)$^\dagger$ & 6.44 (-1\%) \\
grasp 8 mm off-centre sideways & 16 & 7.23--7.53 & 3.41 (-55\%)$^\dagger$ & 5.08 (-32\%)$^\dagger$ \\
grasp 5 mm high on the face & 16 & 7.23--7.53 & 7.17 (-5\%)$^\dagger$ & 7.44 (-1\%) \\
thin box (1 cm tall) & 8 & 5.51--5.74 & 5.35 (-7\%)$^\dagger$ & 6.14 (+7\%)$^\dagger$ \\
cylinder r 2 cm, through the axis & 8 & 5.51--5.74 & 5.30 (-8\%)$^\dagger$ & 5.30 (-8\%)$^\dagger$ \\
cylinder r 2 cm, pad 5 mm above axis & 4 & 4.64--4.84 & 5.34 (+10\%)$^\dagger$ & 5.35 (+11\%)$^\dagger$ \\
cylinder r 2 cm, pad 10 mm above axis & 4 & 4.64--4.84 & 3.53 (-27\%)$^\dagger$ & 3.65 (-25\%)$^\dagger$ \\
mid-air grasp, closing axis level & 16 & 7.23--7.53 & 7.17 (-5\%)$^\dagger$ & 7.41 (-2\%) \\
mid-air, closing axis tilted 15 deg & 16 & 7.48--7.79 & 7.42 (-5\%)$^\dagger$ & 7.64 (-2\%) \\
mid-air, closing axis tilted 30 deg & 16 & 8.36--8.70 & 8.04 (-8\%)$^\dagger$ & 8.31 (-5\%)$^\dagger$ \\
fast lift (10 cm in 0.6 s) & 16 & 6.45--7.53 & 6.52 (-13\%) & 7.02 (-7\%) \\
slow lift (10 cm in 2.4 s) & 16 & 7.45--7.53 & 7.36 (-2\%)$^\dagger$ & 7.44 (-1\%) \\
time step 1 ms & 16 & 7.23--7.53 & 7.17 (-5\%)$^\dagger$ & 7.41 (-2\%) \\
}
}{}
\IfFileExists{../results/calib/characterize_rigid.tex}{% generated by scripts/characterize_contact.py -- do not edit
\newcommand{\calibcharacterizerigid}{%
reference (box 20x4x4 cm) & 16 & 7.23--7.53 & 7.41 (-2\%) & 7.79 (+4\%)$^\dagger$ \\
box yawed 3 deg & 16 & 7.23--7.53 & 1.43 (-81\%)$^\dagger$ & 1.56 (-79\%)$^\dagger$ \\
box yawed 10 deg & 16 & 7.23--7.53 & 1.06 (-86\%)$^\dagger$ & 1.63 (-78\%)$^\dagger$ \\
grasp +1 cm toward the weight & 16 & 8.23--8.53 & 8.34 (-2\%) & never \\
grasp -1 cm away from the weight & 16 & 6.23--6.53 & 6.45 (-1\%) & 6.87 (+5\%)$^\dagger$ \\
grasp 8 mm off-centre sideways & 16 & 7.23--7.53 & 1.06 (-86\%)$^\dagger$ & 1.53 (-80\%)$^\dagger$ \\
grasp 5 mm high on the face & 16 & 7.23--7.53 & 7.42 (-1\%) & 7.84 (+4\%)$^\dagger$ \\
thin box (1 cm tall) & 8 & 5.51--5.74 & 5.84 (+2\%)$^\dagger$ & 6.24 (+9\%)$^\dagger$ \\
cylinder r 2 cm, through the axis & 8 & 5.51--5.74 & 4.13 (-28\%)$^\dagger$ & 7.97 (+39\%)$^\dagger$ \\
cylinder r 2 cm, pad 5 mm above axis & 4 & 4.64--4.84 & 4.72 (-3\%) & 7.57 (+56\%)$^\dagger$ \\
cylinder r 2 cm, pad 10 mm above axis & 4 & 4.64--4.84 & 3.11 (-36\%)$^\dagger$ & 3.11 (-36\%)$^\dagger$ \\
mid-air grasp, closing axis level & 16 & 7.23--7.53 & 7.42 (-1\%) & 7.79 (+4\%)$^\dagger$ \\
mid-air, closing axis tilted 15 deg & 16 & 7.48--7.79 & 7.66 (-2\%) & 8.02 (+3\%)$^\dagger$ \\
mid-air, closing axis tilted 30 deg & 16 & 8.36--8.70 & 8.39 (-4\%) & never \\
fast lift (10 cm in 0.6 s) & 16 & 6.45--7.53 & 6.94 (-8\%) & 7.82 (+4\%)$^\dagger$ \\
slow lift (10 cm in 2.4 s) & 16 & 7.45--7.53 & 7.46 (-1\%) & 7.74 (+3\%)$^\dagger$ \\
time step 1 ms & 16 & 7.23--7.53 & 7.41 (-2\%) & 7.79 (+4\%)$^\dagger$ \\
}
}{}
\IfFileExists{../results/m3/numbers.tex}{% generated by scripts/panda_m3_numbers.py -- do not edit
\newcommand{\mthreeposes}{27}
\newcommand{\mthreelam}{2.18}
\newcommand{\mthreelamlo}{1.79}
\newcommand{\mthreelamhi}{3.42}
\newcommand{\mthreelamnc}{1.23}
\newcommand{\mthreeerrshort}{2.49}
\newcommand{\mthreeerrlong}{1.11}
\newcommand{\mthreeerrlowshort}{0.50}
\newcommand{\mthreeerrhighlong}{2.23}
\newcommand{\mthreeneed}{940}
\newcommand{\mthreeneeds}{1.9}
\newcommand{\mthreedq}{0.051}
\newcommand{\mthreedqlo}{0.042}
\newcommand{\mthreedqhi}{0.111}
\newcommand{\mthreenotare}{0.34}
\newcommand{\mthreetare}{0.015}
\newcommand{\mthreebias}{0.11}
\newcommand{\mthreebiasninety}{0.24}
}{}

\begin{document}

\begin{center}
{\sffamily\bfseries\LARGE AnyGrasp++}\\[3pt]
{\sffamily\large Derivations, step by step}\\[8pt]
Dawit Chun\\[1pt]
{\small Build companion \,|\, supersedes \texttt{proposal/derivations.tex} \,|\, version 2, November 17, 2026}
\end{center}
\vspace{-4pt}
\noindent\rule{\textwidth}{0.4pt}

\textbf{How to use this.} Every equation the code uses, or will use, is derived here with no skipped steps. Each derivation lists its assumptions first. A ``Check'' label points to \texttt{scripts/check\_derivations.py}, which tests the claim independently. Its tags say how:
\begin{itemize}
  \item \texttt{[sym]}: sympy redoes the algebra from scratch;
  \item \texttt{[num]}: Monte Carlo or a different algorithm;
  \item \texttt{[mj]}: read from MuJoCo itself;
  \item \texttt{[arith]}: confirms a quoted number only.
\end{itemize}
``Code'' names the function that implements the result. Shaded boxes are open questions; orange boxes are corrections. Tables of results are generated by the scripts and read in here, so they cannot go stale. Design changes from the proposal, with their reasons, are in \texttt{docs/decisions.md} (D1--D12).

\textbf{What changed since version 1 (after two independent reviews).}
\begin{itemize}
  \item New: the realistic patch pad (A.12) and the exact uniform prior (C.5).
  \item Corrected: load sharing (A.5); the ellipse is an \emph{outer} approximation (A.7); the chatter mechanism (B.4); contact depth (B.5); the CoM accuracy, which was 3$\times$ too optimistic (E.5); the tare numbers (E.2); the statistics under clustering (G.1, G.2); separation in the logistic fit (G.4).
  \item Every check that only repeated its own formula was replaced by an independent test.
\end{itemize}

\tableofcontents

% =====================================================================
\section{Grasp statics and the slip distance}

\subsection{Notation and the two pad models}
World frame with $\ez$ up and gravity $-g\,\ez$, $g = 9.81$\,m/s$^2$. A grasp has centre $\bm p$ (midpoint between the pads) and closing axis $\ah$ (unit vector along which the fingers move). The object has mass $m$ and centre of mass (CoM) $\bm c$; $\bm r = \bm c - \bm p$. Each pad presses with normal force $f_n$ and friction coefficient $\mu$. While the hand accelerates upward at $a_\ell$, the effective weight is $W = m(g+a_\ell)$ (A.6).

Two pad models are used:
\begin{itemize}
  \item \textbf{Soft finger} (A.4--A.10): one contact per pad with a disc-shaped pressure patch of radius $a$. This is the textbook model. In the simulator it is the ``point pad'', used to verify the simulator.
  \item \textbf{Patch pad} (A.12): a set of contact points across the real 17$\times$17\,mm Franka fingertip face, sharing the normal force. This is the ``grid pad'' used for the study (decisions D4).
\end{itemize}
Nominal values: $\mu = 0.5$, $f_n = 70$\,N, $a = 6$\,mm, $m = 0.6$\,kg.

\subsection{The twisting torque and the lever arm $d$}
\textbf{Claim.} Gravity twists the object about the closing axis with torque $\tau_a = -W\,(\bm r\cdot\bm h)$, where $\bm h = \ez\times\ah$. So $|\tau_a| = W|d|$ with $d = \bm r\cdot\bm h$.

\textbf{Assumptions.} Rigid object; the weight acts at the CoM.

\textbf{Derivation.}
\begin{enumerate}
  \item The weight $\bm F = -W\ez$ acts at $\bm c$. Its moment about $\bm p$ is $\bm n = \bm r\times\bm F = -W\,(\bm r\times\ez)$.
  \item The part that turns the object between the pads is the component along the closing axis: $\tau_a = \ah\cdot\bm n = -W\,\ah\cdot(\bm r\times\ez)$.
  \item A scalar triple product is a determinant, $\bm u\cdot(\bm v\times\bm w) = \det[\bm u\ \bm v\ \bm w]$. A cyclic shift of the columns is two swaps, so the sign stays: $\ah\cdot(\bm r\times\ez) = \bm r\cdot(\ez\times\ah)$.
  \item Hence $\tau_a = -W\,\bm r\cdot\bm h$. \hfill $\square$
\end{enumerate}
For a horizontal $\ah$, $\bm h$ is a horizontal unit vector perpendicular to $\ah$. The components of $\bm r$ along $\ah$ and $\ez$ drop out, so $d$ is the horizontal distance from the CoM to the vertical plane containing the closing axis.
\meta{A.2 [num]}{physics.lever\_arm, physics.hinge\_vector}

\subsection{Tilted closing axes}
\textbf{Claim.} If $\ah$ is tilted by $\phi$ above horizontal, then $|\bm h| = \cos\phi$.

\textbf{Derivation.} Write $\ah = \cos\phi\,\bm u + \sin\phi\,\ez$ with $\bm u$ a horizontal unit vector. Then $\bm h = \cos\phi\,(\ez\times\bm u) + \sin\phi\,(\ez\times\ez) = \cos\phi\,(\ez\times\bm u)$. \hfill $\square$

Both the twisting torque and the weight lying in the pad plane scale with $\cos\phi$. So for small tilts the slip distance is that of an object of weight $W\cos\phi$, which the simulator confirms at 15$^\circ$ and 30$^\circ$ (R.3). The weight component $W\sin\phi$ along $\ah$ loads one pad's normal force more than the other's; this is not modelled. The study keeps only grasps whose approach is within 30$^\circ$ of straight down (decisions D3), so $\phi\le30^\circ$.
\meta{A.3 [num]}{physics.hinge\_vector, physics.approach\_ok}

\subsection{Twisting resistance of a disc pad (soft finger)}
\textbf{Claim.} A disc pad of radius $a$ with uniform pressure resists at most $\tau_0 = \tfrac23\mu f_n a$ of twist about its normal. With Hertzian pressure it resists $\tfrac{3\pi}{16}\mu f_n a \approx 0.589\,\mu f_n a$, which is 11.6\% less.

\textbf{Assumptions.} At the start of twisting the object turns about the patch centre, every point slides, and Coulomb friction $\mu\,p\,dA$ acts tangentially everywhere.

\textbf{Derivation (uniform).}
\begin{enumerate}
  \item The pressure is $p_0 = f_n/(\pi a^2)$.
  \item A ring of radius $\rho$ and width $d\rho$ has area $2\pi\rho\,d\rho$ and friction $\mu p_0\,2\pi\rho\,d\rho$ acting at lever arm $\rho$.
  \item Summing the rings: $\tau_0 = \int_0^a 2\pi\mu p_0\rho^2\,d\rho = \tfrac{2\pi}{3}\mu p_0 a^3 = \tfrac23\mu f_n a$. \hfill $\square$
\end{enumerate}
\textbf{Derivation (Hertzian).} $p = p_h\sqrt{1-\rho^2/a^2}$.
\begin{enumerate}
  \item Normalization: with $u = 1-\rho^2/a^2$, $\int_0^a p\,2\pi\rho\,d\rho = 2\pi p_h\frac{a^2}{2}\int_0^1\sqrt u\,du = \frac{2\pi}{3}p_ha^2 = f_n$, so $p_h = \frac{3f_n}{2\pi a^2}$.
  \item Torque: with $\rho = a\sin\theta$, $\int_0^a\rho^2\sqrt{1-\rho^2/a^2}\,d\rho = a^3\int_0^{\pi/2}\sin^2\theta\cos^2\theta\,d\theta = a^3\cdot\frac14\cdot\frac\pi4 = \frac{\pi a^3}{16}$. So $\tau_0 = 2\pi\mu p_h\frac{\pi a^3}{16} = \frac{3\pi}{16}\mu f_na$. \hfill $\square$
\end{enumerate}
\meta{A.4 [sym] (Hertz integrated directly in $\rho$, without the substitution)}{physics.torsional\_mu}

\subsection{How the two pads share the load}
\textbf{Claim.} If the CoM lies in the vertical plane of the closing axis ($\alpha = 0$ in A.11), the best a grasp can do is to have each pad carry $f_t = W/2$ and $\tau = Wd/2$, so the largest holdable lever arm is $\dstar$ (A.8). A CoM offset $\alpha$ \emph{along} the axis lowers it slightly.

\textbf{Derivation.}
\begin{enumerate}
  \item Vertical force balance: $f_{t1} + f_{t2} = W$. Moment balance about $\bm h$ (A.11 step 4): $f_{t1} - f_{t2} = 2W\alpha/w$, with $w$ the object width. So statics \emph{fixes} the force split: $f_{t1,2} = W/2 \pm W\alpha/w$.
  \item Moment balance about $\ah$: the pad centres lie on the axis, so only the pad torques count: $\tau_1 + \tau_2 = Wd$. This split is \emph{not} fixed by statics.
  \item By the ellipse (A.7), a pad carrying $f_t$ can resist at most $\tau_{\max}(f_t) = \tau_0\sqrt{1-(f_t/\mu f_n)^2}$, so the largest holdable $d$ is $[\tau_{\max}(f_{t1}) + \tau_{\max}(f_{t2})]/W$.
  \item $\tau_{\max}$ is concave in $f_t$ (it is a semi-ellipse). For a fixed sum $f_{t1}+f_{t2} = W$, the sum of a concave function is largest at the equal split $\alpha = 0$, where it equals $2\tau_{\max}(W/2) = W\dstar$. \hfill $\square$
\end{enumerate}
\textbf{Numbers} ($w = 4$\,cm): $\dstar = 4.7402,\ 4.7391,\ 4.7360,\ 4.7232$\,cm for $\alpha = 0,\ 0.5,\ 1,\ 2$\,cm.
\fix{Version 1 said ``two equations, four unknowns'' and treated both splits as free. The force split is fixed by moment balance about $\bm h$; only the torque split is free. The conclusion that $\dstar$ is an upper bound stands.}
\meta{A.5 [num] (maximizes over all splits) [arith]}{physics.dstar}

\subsection{Effective weight while lifting}
\textbf{Claim.} While the hand accelerates upward at $a_\ell$, replace $mg$ by $W = m(g+a_\ell)$ everywhere.

\textbf{Derivation.} The object moves with the hand, so $\bm F_{\text{pads}} - mg\ez = ma_\ell\ez$, which gives $\bm F_{\text{pads}} = m(g+a_\ell)\ez$. The extra term $ma_\ell$ acts at the CoM, exactly like extra gravity, so it also adds to the twisting torque. \hfill $\square$

The lift used in step 1 peaks at $a_\ell = 0.40$\,m/s$^2$ (D.1). The measured slip onset should therefore fall between $\dstar$ at that acceleration and $\dstar$ at rest. Faster and slower lifts confirm this (R.3).
\meta{A.6 [arith]; physical test: R.3 fast/slow lift}{physics.dstar(accel=...)}

\subsection{Combined sliding and twisting: the limit surface}
\textbf{Model.} Following Howe and Cutkosky, approximate the loads one disc pad can carry by an ellipse:
\begin{equation}
\left(\frac{f_t}{\mu f_n}\right)^2 + \left(\frac{\tau}{\tau_0}\right)^2 \le 1 .
\label{eq:ellipse}
\end{equation}
Its intercepts are exact: pure sliding at $\mu f_n$, pure twisting at $\tau_0$ (A.4). In between, the exact limit surface of a uniform-pressure disc lies slightly \emph{inside} the ellipse. The ellipse is therefore an \emph{outer}, slightly optimistic approximation: at nominal values it overstates $\dstar$ by 0.18\% (exact 4.7318\,cm vs 4.7402\,cm, computed with A.12 on a fine disc). The exact limit surface of any patch is derived in A.12.
\note{MuJoCo's contact cone \emph{is} this ellipse (B.1). On the point pad the simulator therefore obeys \eqref{eq:ellipse} by construction, and the point-pad tests in R verify the simulator, not the friction model. On the grid pad the twisting resistance instead \emph{emerges} from the spread of the contacts (A.12).}
\fix{Version 1 called the ellipse an approximation without a direction. It overestimates the combined capacity by up to about 5\% in radius; the effect on $\dstar$ at nominal values is 0.18\%.}
\meta{A.7 [num] (exact disc via A.12 vs ellipse; pure twist = $\frac23\mu f_na$)}{--}

\subsection{Solving for the slip distance $\dstar$ (soft finger)}
\textbf{Claim.}
\begin{equation}
\dstar = \frac{4\mu f_n a}{3W}\sqrt{1-\left(\frac{W}{2\mu f_n}\right)^2},\qquad\text{which requires } W < 2\mu f_n .
\label{eq:dstar}
\end{equation}
\textbf{Derivation.}
\begin{enumerate}
  \item Put $f_t = W/2$ and $\tau = Wd/2$ (A.5) into \eqref{eq:ellipse}: $\left(\frac{W}{2\mu f_n}\right)^2 + \left(\frac{Wd/2}{\frac23\mu f_n a}\right)^2\le1$.
  \item Simplify the second fraction: $\frac{Wd}{2}\cdot\frac{3}{2\mu f_na} = \frac{3Wd}{4\mu f_na}$.
  \item The remaining room, $1-\left(\frac{W}{2\mu f_n}\right)^2$, must be non-negative, otherwise even $d = 0$ fails: $W\le2\mu f_n$, i.e.\ $m\le7.14$\,kg at rest.
  \item Take square roots and solve for $d$. \hfill $\square$
\end{enumerate}
\textbf{Numbers.} $W = 5.886$\,N gives a prefactor of 4.757\,cm and a root of 0.99646, so $\dstar = 4.74$\,cm. The simple form (without the root) gives 4.76\,cm, 0.35\% higher.

\begin{table}[H]
\centering
\caption{Soft-finger $\dstar$ in cm, $f_n = 70$\,N, $a = 6$\,mm, at rest}
\footnotesize
\begin{tabular}{@{}lccc@{}}
\toprule
$m$ [kg] & $\mu = 0.3$ & $\mu = 0.5$ & $\mu = 0.8$ \\
\midrule
0.3 & 5.69 & 9.51 & 15.22 \\
0.6 & 2.83 & \textbf{4.74} & 7.60 \\
1.0 & 1.67 & 2.83 & 4.55 \\
\bottomrule
\end{tabular}
\end{table}
\meta{A.8 [sym] [arith]}{physics.dstar}

\subsection{What $\dstar$ depends on}
\textbf{Claim.}
\begin{enumerate}[label=(\roman*)]
  \item $\dstar$ is exactly proportional to $a$.
  \item It depends on $\mu$ and $f_n$ only through $k = \mu f_n$, with elasticity $\partial\ln\dstar/\partial\ln k = 1/(1-x^2)$, $x = W/(2k)$. At nominal values this is 1.007.
  \item It falls roughly as $1/m$.
\end{enumerate}
\textbf{Derivation.} Write $\dstar = \frac{4a}{3W}k\sqrt{1-x^2}$, so $\ln\dstar = \text{const}+\ln k+\tfrac12\ln(1-x^2)$. Since $d\ln x/d\ln k = -1$, the last term contributes $\tfrac12\cdot\frac{-2x^2}{1-x^2}\cdot(-1) = \frac{x^2}{1-x^2}$. Adding the 1 from $\ln k$ gives the claim. \hfill $\square$
\meta{A.9 [num] (finite-difference elasticity; $\dstar(0.25, 140) = \dstar(0.5, 70)$)}{physics.dstar}

\subsection{The grip force a grasp needs}
\textbf{Claim.} To hold lever arm $d$ with a soft-finger pad, each pad needs at least $f_{n,\min} = \frac{W}{\mu}\sqrt{\frac14+\left(\frac{3d}{4a}\right)^2}$.

\textbf{Derivation.} Take the boundary in A.8 step 1. Factor out $(W/\mu f_n)^2$: $\left(\frac{W}{\mu f_n}\right)^2\left[\frac14 + \left(\frac{3d}{4a}\right)^2\right] = 1$. Solve for $f_n$. \hfill $\square$

At $d = 0$ this gives $W/(2\mu)$, i.e.\ each pad carries half the weight by friction. For $d\gg4$\,mm it is about $3Wd/(4\mu a)$. Example: $d = 4$\,cm needs 59.2\,N.
\note{The M3 test lift itself must use the full 70\,N. At 30\,N the slip distance is only 2\,cm, so the object would twist \emph{during} the measurement (decisions D8). This formula is for lowering the grip \emph{after} the measurement.}
\meta{A.10 [sym] [num] (round trip)}{physics.required\_grip}

\subsection{Why twisting is the weak mode}
\textbf{Claim.} A CoM offset $\alpha$ along the closing axis tilts the object in the plane of the fingers. The pads resist up to $(\mu f_n - W/2)\,w$ of that moment, which is 1.28\,N\,m for $w = 4$\,cm, against 0.28\,N\,m for twisting.

\textbf{Derivation.}
\begin{enumerate}
  \item Write $\bm r_h = \alpha\ah + \beta\bm h$. Then $\bm n = W(\ez\times\bm r)$.
  \item $\ez\times\ah = \bm h$, and $\ez\times\bm h = -\ah$ (BAC--CAB with $\ez\cdot\ah = 0$). So $\bm n = W(\alpha\bm h - \beta\ah)$.
  \item The $\bm h$ part must be balanced by vertical pad forces at $\pm\frac w2\ah$: $f_1+f_2 = W$ and $f_1 - f_2 = 2W\alpha/w$.
  \item Coulomb on the more loaded pad: $W/2 + W\alpha/w\le\mu f_n$. \hfill $\square$
\end{enumerate}
\meta{A.11 [num] (vector identity) [arith]}{(design argument)}

\subsection{Patch pads: twisting resistance from the spread of the contacts}
This is the model behind the realistic grid pad. The pad touches the object at points $\bm p_k$ (2-D, in the pad plane: $x$ along the object, $z$ up, origin at the pad centre). Point $k$ carries normal load $n_k$ (equal shares $f_n/K$ by default) and Coulomb friction $|\bm f_k|\le\mu n_k$.

\textbf{Claim 1 (exact limit).} The largest twisting moment one pad can resist while also pushing up with force $f_{\text{up}}$ is
\begin{equation}
T_{\max}(f_{\text{up}}) = \min_{\bm c\in\R^2}\ \Big[\sum_k \mu n_k\,|\bm p_k - \bm c| + c_x f_{\text{up}}\Big],
\label{eq:dual}
\end{equation}
and the minimizing $\bm c$ is the point the object turns about when it starts to slip. With equal sharing between the two pads (optimal by A.5), $\dstar = 2T_{\max}(W/2)/W$.

\textbf{Claim 2 (pure twist).} $T_{\max}(0) = \mu\sum_k n_k|\bm p_k - \bm c_{\text{med}}|$, where $\bm c_{\text{med}}$ is the weighted geometric median. For a symmetric patch this is $\mu f_n\cdot$(mean distance to the centre).

\textbf{Derivation.}
\begin{enumerate}
  \item Each contact force $\bm f_k$ contributes the vertical/horizontal force $\bm f_k$ and the moment $p_{kx}f_{kz} - p_{kz}f_{kx}$ about the pad centre. Write this as $L_k\bm f_k$, with $L_k$ a $3\times2$ matrix. The set $S$ of loads $(F_x, F_z, T)$ the pad can carry is the sum of the images of the discs $\{|\bm f_k|\le\mu n_k\}$. It is convex and compact.
  \item Support function: for any direction $\bm u$, $\max_{\bm s\in S}\bm u\cdot\bm s = \sum_k\max_{|\bm f|\le\mu n_k}(L_k^{\mathsf T}\bm u)\cdot\bm f = \sum_k\mu n_k|L_k^{\mathsf T}\bm u|$.
  \item We want $\max\{T : (0, f_{\text{up}}, T)\in S\}$. Lagrange duality for maximizing a linear function over a convex compact set with two linear equalities gives $\min_{\lambda\in\R^2}\big[\max_{\bm s\in S}(T + \lambda_xF_x + \lambda_zF_z) - \lambda_zf_{\text{up}}\big]$. Strong duality holds because $f_{\text{up}} < \mu f_n$ puts the target force in the interior of the achievable forces.
  \item With $\bm u = (\lambda_x, \lambda_z, 1)$: $L_k^{\mathsf T}\bm u = (\lambda_x - p_{kz},\ \lambda_z + p_{kx})$. Substitute $\bm c = (-\lambda_z, \lambda_x)$: this is $(c_z - p_{kz},\ p_{kx} - c_x)$, whose length is $|\bm p_k - \bm c|$. Also $-\lambda_zf_{\text{up}} = c_xf_{\text{up}}$, which gives \eqref{eq:dual}.
  \item Pure twist, $f_{\text{up}} = 0$: \eqref{eq:dual} is the weighted geometric median problem. For a patch symmetric under $\bm p\to-\bm p$ the median is the centre. \hfill $\square$
\end{enumerate}
\textbf{Interpretation.} The object starts to slip by turning about $\bm c$; each point's friction opposes its own sliding direction; \eqref{eq:dual} picks the $\bm c$ that delivers exactly the needed lift.

\textbf{Numbers} (4$\times$4 grid over the 17$\times$17\,mm face, $\mu = 0.5$, 70\,N, 0.6\,kg):

\begin{tabular}{@{}ll@{}}
Full face (16 points) & $\dstar = 7.53$\,cm \\
Two middle rows, a line contact on a handle (8 points) & 5.74\,cm \\
One row, off-centre (4 points) & 4.84\,cm \\
\end{tabular}

The grid's mean point distance is 6.37\,mm, against 6.51\,mm for continuous uniform pressure on the square: a 2\% discretization error. MuJoCo shares the force equally over the 16 points (4.37--4.38\,N each).
\note{After slipping starts the object turns by $\theta$, and its lever arm shrinks to about $d\cos\theta$. The grid-pad object then re-sticks once $d\cos\theta$ falls back below $\dstar$. That is why its 10$^\circ$ failure point lies 4--10\% beyond its slip onset, while the point pad overshoots at once (R.1, R.2).}
\meta{A.12 [num] (duality formula vs a linear program and vs centre-of-rotation root finding; pure twist; continuous square) [mj] (equal sharing)}{physics.patch\_torque\_capacity, physics.patch\_dstar}

% =====================================================================
\section{The simulator model}

\subsection{MuJoCo's friction cone, and which geom's parameters a contact uses}
\textbf{Claim 1.} With \texttt{condim="4"} and \texttt{cone="elliptic"}, MuJoCo's contact obeys \eqref{eq:ellipse} exactly when $\mu_1 = \mu_2 = \mu$ and the torsional coefficient is $\mu_3 = \tfrac23\mu a$.

\textbf{Derivation.} The elliptic cone requires $f_n^2\ge(f_1/\mu_1)^2 + (f_2/\mu_2)^2 + (\tau/\mu_3)^2$ (MuJoCo documentation, \emph{Computation}). With $\mu_1 = \mu_2 = \mu$, divide by $f_n^2$ and compare with \eqref{eq:ellipse}. \hfill $\square$

\textbf{Claim 2.} When two geoms touch with equal \texttt{priority}, MuJoCo uses the \emph{element-wise maximum} of their friction and the \emph{solmix-weighted average} of their solref/solimp. When the priorities differ, the higher-priority geom supplies all contact parameters.

\textbf{Consequence.} In version 1 the stiff pad settings were averaged with the box's defaults (the pad contact ran at solref 0.012, not 0.004), and an object could never be less slippery than the pad. Now every object geom has priority 1 and owns all its contact parameters (decisions D5).
\meta{B.1 [mj] (the max/average/priority rules; a sliding box sits exactly on the cone, $|f_t|/f_n = \mu$; the rig's pad contacts use the object's values)}{contact\_rig.build\_xml}

\subsection{Placing the hidden weight}
\textbf{Claim.} An invisible point mass $m_h$ at distance $x_h$ from the shell's centroid, with the shell at $m - m_h$, moves the CoM by $d$ when $m_h = md/x_h$.

\textbf{Derivation.} $\frac{(m-m_h)\cdot0 + m_hx_h}{m} = d$. This needs $d<x_h$. \hfill $\square$

An 8\,cm shift needs $x_h\ge8.9$\,cm, so the object must be at least about 18\,cm long.
\note{This point mass is used only by the calibration rig. The study's scanned objects use a hidden \emph{core} of real volume and at most tungsten's density (section H).}
\fix{Version 1 also required the visible shell to keep $\ge10\%$ of the mass. That rule came from the averaged, softer contact (B.1, B.5) and no longer applies.}
\meta{B.2 [mj]}{contact\_rig.build\_xml}

\subsection{The gripper: how a tendon force becomes pad force}
\textbf{Claims.}
\begin{enumerate}[label=(\roman*)]
  \item A tendon force $F$ on $L = \frac12(q_1+q_2)$ pushes each finger with $\frac12F$, so $f_n$ per pad needs $|F| = 2f_n$.
  \item Menagerie's stock gripper squeezes a 4\,cm object with about 1\,N per finger.
  \item The tendon-space finger inertia is $m_t = 0.23$\,kg.
\end{enumerate}
\textbf{Derivation.}
\begin{enumerate}
  \item Virtual work: $F\,\delta L = F\cdot\frac12(\delta q_1+\delta q_2)$, so each joint gets $\frac12F$.
  \item Menagerie's actuator is $F = 0.0157u - 100L - 10\dot L$. Commanded closed ($u = 0$) and at rest, $F = -100L$. On a 4\,cm box, $L\approx20$\,mm, so each finger pushes $\frac12\cdot100\cdot0.020 = 1.0$\,N.
  \item With $m_f = 0.015 + 0.1$ (armature) $= 0.115$\,kg per finger, $m_f\ddot L = \frac12F$, so $m_t = 2m_f$. Equivalently $1/(J_tM^{-1}J_t^{\mathsf T})$ with $J_t = [\frac12,\frac12]$. \hfill $\square$
\end{enumerate}
\textbf{The force grip used instead:} $F = \mathrm{clip}(k_p(u-L) - k_v\dot L,\ \pm2f_n)$. The target $u$ runs closed at 0.1\,m/s to $-1$\,cm. With $k_p = 2f_n/0.005 = 28{,}000$\,N/m, the clip is active 5\,mm past contact, so the pads get exactly $f_n$. Critical damping gives $k_v = 2\sqrt{k_pm_t} = 160.5$\,N\,s/m.
\meta{B.3 [mj] (inertia from MuJoCo's mass matrix; stock grip force read from MuJoCo: 1.0\,N; force grip 70.0\,N)}{contact\_rig.\_grip\_actuator}

\subsection{Why the first force controller chattered}
\textbf{Claim.} Under MuJoCo's \texttt{implicitfast} integrator, an \emph{unclamped} actuator damper is treated implicitly and is stable for any $k_v$. But ``a saturated forcerange \ldots\ contributes no implicit damping, since the clamped force does not depend on velocity'' (MuJoCo documentation, \emph{Modeling, Actuators}). A clamp-to-clamp oscillation therefore exists if and only if $k_v\Delta t/m\ge2$.

\textbf{Derivation.} Suppose the force sits at a clamp, $\pm F_{\max}$, and alternates. One step changes the velocity by $F_{\max}\Delta t/m$, so the velocity alternates between $\pm v$ with $v = F_{\max}\Delta t/(2m)$. The unclamped force would be $-k_vv$ plus the command. The cycle persists if that keeps exceeding the clamp, i.e.\ $k_vF_{\max}\Delta t/(2m)\ge F_{\max}$, which is $k_v\Delta t/m\ge2$. \hfill $\square$

The first controller had $k_v = 1400$: $1400\cdot0.002/0.23 = 12.2$, which chatters. The current one has $160.5\cdot0.002/0.23 = 1.40$, which does not.
\fix{Version 1 said the velocity ``grows 11-fold per step'' (the explicit-Euler picture) and that the clamp mechanism was unconfirmed. The motion is a bounded two-step cycle, and the mechanism is documented.}
\meta{B.4 [mj] (old controller flips sign every step, new one never) [arith]}{contact\_rig.\_grip\_actuator}

\subsection{How deep a contact sinks}
\textbf{Claim.} At rest, a contact with normal force $f$ sinks to the depth $r$ that solves
\begin{equation}
r = \frac{(1-d(r))\,\hat A\,f\,d_{\max}^2\,t_c^2\,\zeta^2}{d(r)^2},
\end{equation}
where:
\begin{itemize}
  \item $(t_c,\zeta)$ is the contact's solref and $d(r)$ its solimp impedance at depth $r$;
  \item $\hat A$ is MuJoCo's approximation of the contact's inverse inertia: a body-based estimate frozen at the initial pose, \emph{not} the exact $JM^{-1}J^{\mathsf T}$.
\end{itemize}
\textbf{Derivation.}
\begin{enumerate}
  \item MuJoCo's soft constraint solves $(A+R)f = a_{\text{ref}} - a_0$. Here $R = \frac{1-d}{d}\hat A$, and $a_{\text{ref}} = -b\,v - k\,d\,r$ with $k = 1/(d_{\max}^2t_c^2\zeta^2)$ (documentation, \emph{Computation}).
  \item At rest the constraint acceleration $a_0 + Af$ is zero, and $v = 0$. So $Rf = a_{\text{ref}} = -k\,d\,r$.
  \item Substitute $R$ and $k$, and take magnitudes. \hfill $\square$
\end{enumerate}
\textbf{Consequences.} Depth grows with $\hat A$. A heavy weight far from the body's origin raises the body's $\hat A$, which is why version 1 saw depths of 1.2--7.5\,mm grow with the shift. With object-owned contacts, the step-1 grid pad sank 0.11\,mm (rubber) or 0.003--0.007\,mm (rigid). The nominal pad is now the rubber layer of B.6: 0.161\,mm on a flat face of a free 1\,kg object, the physical estimate for a 2\,mm layer (decisions D4, D20).
\fix{Version 1 attributed the trend to the exact effective mass at the contact and could not reproduce the depths (it used the pad's own parameters, not the averaged ones). The formula above reproduces MuJoCo's depth to within 2\%.}
With a direct stiffness (solref $= (-K, -b)$, used from B.6 on), MuJoCo sets $k = K/d_{\max}^2$, so $r = (1-d)\hat A f\,d_{\max}^2/(K\,d^2)$. With a constant impedance $d = d_{\max}$ each contact is a linear spring of stiffness $K/((1-d)\hat A)$.
\fix{The first form of this line omitted MuJoCo's division by $d_{\max}^2$ and was 11\% off; check B.5 caught it. The pad's stiffness was always set by measurement (B.6), so it was not affected.}
\meta{B.5 [mj] (formula vs MuJoCo's depth, for the layer pad, the step-1 pad and default contacts)}{(diagnostic)}

\subsection{The fingertip as a thin rubber layer}
\textbf{Model.} A rubber layer of thickness $h$ and Young's modulus $E$ on a rigid finger is approximated, when the loaded region is wider than $h$, as a bed of independent springs (Winkler foundation; it ignores the shear coupling between neighbouring cells and the stiffening of a bonded, nearly incompressible layer): the pressure is $E\,\delta/h$ where it sinks by $\delta$. One cell of the $4\times4$ grid has area $A_c = (17\text{ mm}/4)^2 = 18.1$\,mm$^2$, so it is a spring
\begin{equation}\label{eq:winkler}
k_c = \frac{E\,A_c}{h} = 27.1\ \text{kN/m}\quad(E = 3\text{ MPa},\ h = 2\text{ mm}).
\end{equation}
Against a surface whose gap to the pad face is $s_i$ at cell $i$, the cells sink $\delta_i = \max(0, \delta_0 - s_i)$, with $\delta_0$ fixed by the squeeze force $F = k_c\sum_i\delta_i$: flat face $s_i = 0$ (each cell $F/16$, sink 0.161\,mm at 70\,N); a face misaligned by $\theta$ across the pad $s_i = x_i\tan\theta$; a cylinder of radius $R$ along the pad $s_i = z_i^2/2R$.

\textbf{In MuJoCo.} By B.5, a constant impedance and a direct stiffness $K$ make every pad point a linear spring of stiffness $K/((1-d)\hat A)$. $\hat A$ is dominated by the finger, so it is nearly the same for every point. $k$ is set so that a free 1\,kg object, the study's condition, sees exactly $k_c$ per point.

\textbf{Why.} The step-1 pad was calibrated only on perfectly aligned flat faces. At 60\,kN/m per point, and stiffening with depth, any misalignment above about $0.4^\circ$ loaded only the edge rows, and the full B0 pipeline ran on 1--2 contact points per pad (decisions D19). Both pads match the Winkler model with their own stiffness; the layer pad has the physical one.
\meta{B.6 [mj] (per-point stiffness $= k_c$ on a free 1\,kg object; misaligned 1/2/4$^\circ$ and a 15\,mm cylinder: engaged points, forces and $\dstar$ match the Winkler solution)}{scripts/calibrate\_pad\_layer.py, contact\_rig.CONTACT}

% =====================================================================
\section{Scoring grasps under CoM uncertainty (M2)}

\subsection{The CoM prior from the point cloud}
\textbf{Claim.} If the CoM is uniform over the object's footprint (an $L\times W_o$ rectangle with axes $V$), its covariance is $V\mathrm{diag}(L^2/12, W_o^2/12)V^{\mathsf T}$.

\textbf{Derivation.} For $x\sim U[-L/2,L/2]$: $\mathrm{Var}\,x = \frac1L\int_{-L/2}^{L/2}x^2dx = \frac{L^2}{12}$. The two axes are independent, and $\mathrm{Cov}(V\bm x') = V\,\mathrm{Cov}(\bm x')V^{\mathsf T}$. \hfill $\square$
\meta{C.1 [sym]}{(M2)}

\subsection{The lever arm under a Gaussian belief}
\textbf{Claim.} If $\bm c\sim\mathcal N(\bm c_0,\Sigma)$, then $d_i\sim\mathcal N\big((\bm c_0-\bm p_i)\cdot\bm h_i,\ \bm h_i^{\mathsf T}\Sigma\bm h_i\big)$.

\textbf{Derivation.} $d_i$ is affine in $\bm c$, and affine functions of Gaussians are Gaussian with mean $\bm h_i^{\mathsf T}\bm c_0 - \bm h_i^{\mathsf T}\bm p_i$ and variance $\bm h_i^{\mathsf T}\Sigma\bm h_i$. \hfill $\square$

This is used \emph{after} the test lift, when the belief is Gaussian (F.1). Before the lift, the prior is uniform; use C.5, not a Gaussian with matched mean and covariance.
\meta{C.2 [num]}{physics.p\_hold}

\subsection{Probability that a grasp holds (Gaussian belief)}
$P_{\text{hold}} = \Pr(|d|\le\dstar) = \Phi\!\left(\frac{\dstar-\mu_d}{\sigma_d}\right) - \Phi\!\left(\frac{-\dstar-\mu_d}{\sigma_d}\right)$. The M2 score is $s_{\text{det}}\cdot P_{\text{hold}}$ (decisions D9). $\dstar$ may be the soft-finger value or a patch value (A.12). $m$, $\mu$ and the patch are treated as known.
\meta{C.3 [num]}{physics.p\_hold}

\subsection{The deterministic margin}
$\rho = 1-|\mu_d|/\dstar$, set to $-\infty$ when $\dstar = 0$, i.e.\ when nothing can hold. F.2 relates it to $P_{\text{hold}}$.
\meta{C.4 [num]}{physics.margin}

\subsection{Exact probability under the uniform footprint prior}
\textbf{Claim.} If $\bm c$ is uniform on the rectangle, then $d = \mu_d + X + Y$ with $X\sim U[-\alpha,\alpha]$ and $Y\sim U[-\beta,\beta]$ independent, where $\alpha = |\bm h\cdot\bm v_1|L/2$ and $\beta = |\bm h\cdot\bm v_2|W_o/2$. The CDF of $X+Y$ is
\begin{equation}
F(t) = \frac{\rho^2(t+\alpha+\beta) - \rho^2(t-\alpha+\beta) - \rho^2(t+\alpha-\beta) + \rho^2(t-\alpha-\beta)}{8\alpha\beta},\qquad\rho(u) = \max(u,0),
\end{equation}
and $P_{\text{hold}} = F(\dstar-\mu_d) - F(-\dstar-\mu_d)$.

\textbf{Derivation.}
\begin{enumerate}
  \item $F(t)$ is the area of $\{(x,y)\in[-\alpha,\alpha]\times[-\beta,\beta] : x+y\le t\}$ divided by $4\alpha\beta$.
  \item The region $\{x\ge x_0,\ y\ge y_0,\ x+y\le t\}$ is a right triangle of area $\frac12\rho(t-x_0-y_0)^2$.
  \item The rectangle is $\mathbf 1[x\ge-\alpha]\mathbf 1[y\ge-\beta] - \mathbf 1[x\ge\alpha]\mathbf 1[y\ge-\beta] - \mathbf 1[x\ge-\alpha]\mathbf 1[y\ge\beta] + \mathbf 1[x\ge\alpha]\mathbf 1[y\ge\beta]$ (inclusion--exclusion). Apply step 2 to each term. \hfill $\square$
\end{enumerate}
If one width is zero, $F$ becomes the uniform CDF of the other; if both are zero, a step.
\fix{Version 1 used a Gaussian with the uniform's mean and covariance without saying so. That misestimates $P_{\text{hold}}$ by up to about 0.11, for example 0.47 vs 0.59 for a 20\,cm footprint.}
\meta{C.5 [num] (vs Monte Carlo; the Gaussian values are printed for comparison)}{physics.p\_hold\_uniform}

% =====================================================================
\section{Motion}

\subsection{The minimum-jerk lift and its peak acceleration}
\textbf{Claims.} $s(\tau) = 10\tau^3 - 15\tau^4 + 6\tau^5$ minimizes $\int(s''')^2$ for a rest-to-rest move. The peak acceleration is $\frac{10\sqrt3}{3}\frac h{T^2}$, at $\tau = \frac{3-\sqrt3}6$.

\textbf{Derivation.}
\begin{enumerate}
  \item The Euler--Lagrange equation for $\int(s''')^2$ is $-\frac{d^3}{d\tau^3}(2s''') = 0$, so $s^{(6)} = 0$ and $s$ is a quintic.
  \item Rest at $\tau = 0$ gives $c_0 = c_1 = c_2 = 0$. Rest at $\tau = 1$ gives $c_3+c_4+c_5 = 1$, $3c_3+4c_4+5c_5 = 0$ and $6c_3+12c_4+20c_5 = 0$. These solve to $c_3 = 10$, $c_4 = -15$, $c_5 = 6$.
  \item $s'' = 60\tau(1-\tau)(1-2\tau)$. Setting $s''' = 0$ gives $\tau^2-\tau+\frac16 = 0$. At the smaller root, $\tau(1-\tau) = \frac16$ and $1-2\tau = \frac{\sqrt3}3$, so $s'' = \frac{10\sqrt3}3$. \hfill $\square$
\end{enumerate}
For 10\,cm in 1.2\,s this is 0.401\,m/s$^2$.
\meta{D.1 [sym]; [num] (50 random admissible perturbations all have more jerk)}{rig\_trial.min\_jerk, a\_peak}

\subsection{Damped least-squares IK (M1)}
\textbf{Claim.} The step minimizing $\|J\Delta\bm q-\bm e\|^2+\lambda^2\|\Delta\bm q\|^2$ is $\Delta\bm q = J^{\mathsf T}(JJ^{\mathsf T}+\lambda^2I)^{-1}\bm e$.

\textbf{Derivation.} Setting the gradient to zero gives $(J^{\mathsf T}J+\lambda^2I)\Delta\bm q = J^{\mathsf T}\bm e$. Then use the push-through identity $J^{\mathsf T}(JJ^{\mathsf T}+\lambda^2I) = (J^{\mathsf T}J+\lambda^2I)J^{\mathsf T}$. \hfill $\square$

Manipulability: $\sqrt{\det J_vJ_v^{\mathsf T}} = \prod\sigma_i$.
\meta{D.2 [num]}{scripts/panda\_m3\_numbers.ik}

\subsection{How far the object turned in the hand}
\textbf{Claim.} A rotation matrix $R$ turns by $\theta = \arccos\big((\tr R-1)/2\big)$.

\textbf{Derivation.} By Rodrigues, $\tr R = \tr I + \sin\theta\,\tr[\bm k]_\times + (1-\cos\theta)\tr[\bm k]_\times^2 = 3 + 0 - 2(1-\cos\theta) = 1 + 2\cos\theta$. \hfill $\square$

\textbf{What the trial measures.} It uses $Q = R_{\text{rel}}R_{\text{rel},0}^{\mathsf T}$, with $R_{\text{rel}} = R_{\text{hand}}^{\mathsf T}R_{\text{obj}}$. This is the object's rotation relative to the hand since the load came on, expressed in the hand frame. Its rotation vector $\bm\theta = \theta\bm k$ splits in two:
\begin{itemize}
  \item \emph{tilt} $\sqrt{\theta_x^2+\theta_y^2}$: twist about the closing axis $y$, plus tilt in the finger plane;
  \item \emph{squaring} $|\theta_z|$: rotation about the approach axis $z$, i.e.\ the object lining up with the fingers, which is harmless.
\end{itemize}
Failure is a tilt above 10$^\circ$ (decisions D2).
\meta{D.3 [num]}{rig\_trial.run\_trial (rot\_in\_hand)}

% =====================================================================
\section{The test lift (M3)}

\subsection{Arm statics with a held object}
\textbf{Claim.} At rest, $\bm\tau = \bm g(\bm q) - J^{\mathsf T}\bm w$, where $\bm w = [\bm f;\bm n]$ is the wrench the object applies at $\bm p$.

\textbf{Derivation.} Dynamics: $M\ddot{\bm q}+C\dot{\bm q}+\bm g = \bm\tau + J^{\mathsf T}\bm w$. The last term follows from virtual work: $\bm w^{\mathsf T}J\delta\bm q$. Set $\dot{\bm q} = \ddot{\bm q} = 0$. \hfill $\square$
\meta{E.1 [mj] (MuJoCo inverse dynamics with an applied wrench at the grasp point)}{(M3)}

\subsection{Zeroing the scale (tare)}
\textbf{Claim.} $\bm y = [\hat{\bm g}(\bm q_1)-\hat{\bm g}(\bm q_0)] - [\bm\tau^{(1)}-\bm\tau^{(0)}] = J_1^{\mathsf T}\bm w + [\delta\bm g(\bm q_1)-\delta\bm g(\bm q_0)] - (\bm\epsilon_1-\bm\epsilon_0)$. The model error enters only as a difference, which is first order in $\Delta\bm q$.

\textbf{Derivation.}
\begin{enumerate}
  \item Readings: $\bm\tau^{(k)} = \bm g(\bm q_k) - J_k^{\mathsf T}\bm w_k + \bm\epsilon_k$, with $\bm w_0 = 0$ (object still on the table) and $\bm w_1 = \bm w$.
  \item The robot's model is $\hat{\bm g} = \bm g+\delta\bm g$.
  \item Subtract: the true gravity terms cancel. Taylor-expand $\delta\bm g(\bm q_1)-\delta\bm g(\bm q_0)\approx\frac{\partial\delta\bm g}{\partial\bm q}\Delta\bm q$. \hfill $\square$
\end{enumerate}
\textbf{Numbers (real Panda, \mthreeposes\ lift poses).}
\begin{itemize}
  \item A 2\,cm lift moves the largest joint by \mthreedq\,rad (median; range \mthreedqlo--\mthreedqhi).
  \item With random $\pm5\%$ link-mass errors, the gravity-model error is \mthreenotare\,N\,m without the tare and \mthreetare\,N\,m with it (medians).
  \item The CoM bias left over is \mthreebias\,cm (median), \mthreebiasninety\,cm (90th percentile).
\end{itemize}
\fix{Version 1 said the joints move 0.01--0.02\,rad, and its bias numbers came from a made-up gravity model.}
\meta{E.2 [num] (from scripts/panda\_m3\_numbers.py)}{(M3)}

\subsection{Least-squares wrench and its covariance}
\textbf{Claim.} $\hat{\bm w} = (JJ^{\mathsf T})^{-1}J\bar{\bm y}$, with $\mathrm{Cov}(\hat{\bm w}) = \frac{2\sigma_\tau^2}N(JJ^{\mathsf T})^{-1}$. It is unbiased except for the tare residual $P[\delta\bm g(\bm q_1)-\delta\bm g(\bm q_0)]$.

\textbf{Derivation.}
\begin{enumerate}
  \item Minimizing $\|J^{\mathsf T}\bm w-\bar{\bm y}\|^2$ gives $JJ^{\mathsf T}\bm w = J\bar{\bm y}$, so $\hat{\bm w} = P\bar{\bm y}$ with $P = (JJ^{\mathsf T})^{-1}J$, and $PJ^{\mathsf T} = I$.
  \item Each hold averages $N$ white samples, giving variance $\sigma_\tau^2/N$. A difference of two such averages has $2\sigma_\tau^2/N$.
  \item So the covariance is $P\frac{2\sigma_\tau^2}NP^{\mathsf T} = \frac{2\sigma_\tau^2}N(JJ^{\mathsf T})^{-1}$. \hfill $\square$
\end{enumerate}
\meta{E.3 [num] (simulated per sample: two holds of $N$ samples, averaged and differenced)}{(M3)}

\subsection{Mass and CoM from the wrench}
\textbf{Claims.}
\begin{itemize}
  \item $\hat m = -\hat f_z/g$.
  \item $\bm r_\perp = \frac{\bm f\times\bm n}{\|\bm f\|^2}$, which for $\bm f = -mg\ez$ is $\frac{\bm n\times\ez}{mg}$.
  \item The CoM height is unobservable from one pose.
  \item $\bm n\cdot\bm f = 0$ serves as a consistency check.
\end{itemize}

\textbf{Derivation.}
\begin{enumerate}
  \item BAC--CAB: $\bm f\times(\bm r\times\bm f) = \bm r\|\bm f\|^2 - \bm f(\bm f\cdot\bm r) = \|\bm f\|^2\bm r_\perp$.
  \item For $\bm f = -mg\ez$: $\bm f\times\bm n = mg(\bm n\times\ez)$.
  \item Example: $\bm r = (1,0,0)$ gives $\bm n = (0,mg,0)$, and $(\bm n\times\ez)/(mg) = (1,0,0)$, which is correct.
  \item $\bm r\to\bm r+\lambda\bm f$ leaves $\bm n$ unchanged, so the height cannot be seen.
  \item $\bm n\cdot\bm f = (\bm r\times\bm f)\cdot\bm f = 0$. \hfill $\square$
\end{enumerate}
\fix{\texttt{proposal/derivations.tex} Sec.~4.4 wrote $(\ez\times\bm n)/(mg)$, which has the wrong sign: it would place the CoM on the wrong side of the grasp.}
\meta{E.4 [num] (includes the sign case)}{(M3)}

\subsection{Expected CoM accuracy}
\textbf{Claim.} Per horizontal axis, $\sigma_{r,i} = \frac{\sigma_\tau\sqrt{2/N}\sqrt{\lambda_i}}{mg}$, where $\lambda_i$ is the matching diagonal entry of the angular block of $(JJ^{\mathsf T})^{-1}$ at the lift pose.

\textbf{Derivation.} $\bm r_\perp = (\bm n\times\ez)/(mg)$ rotates the horizontal part of $\bm n$ by $90^\circ$ and scales it by $1/(mg)$. The covariance of $\bm n$ is the angular block of E.3. \hfill $\square$

\textbf{Why version 1 was 3$\times$ too optimistic.} It replaced $\lambda$ by $3/7$ ($\sqrt{3/7} = 0.65$). That number assumed the joint axes point in random directions, and it ignored the coupling between the linear and angular rows of $J$.
\begin{itemize}
  \item By the Schur complement, the coupled angular block $(J_\omega(I-P_v)J_\omega^{\mathsf T})^{-1}$ is always at least the uncoupled $(J_\omega J_\omega^{\mathsf T})^{-1}$.
  \item Averaging the axes before inverting (Jensen's inequality) underestimates the inverse too.
\end{itemize}

\textbf{Real Panda, \mthreeposes\ top-down lift poses.} The median $\sqrt\lambda$ is \mthreelam\ (range \mthreelamlo--\mthreelamhi; \mthreelamnc\ without coupling), against the guessed 0.65. The median 2-D CoM error is:
\begin{itemize}
  \item \mthreeerrshort\,cm at $\sigma_\tau = 0.5$\,N\,m with a 0.3\,s hold;
  \item \mthreeerrlong\,cm with a 1.5\,s hold;
  \item \mthreeerrlowshort\,cm at $\sigma_\tau = 0.1$\,N\,m with 0.3\,s.
\end{itemize}
Getting below 1\,cm at 0.5\,N\,m needs about \mthreeneed\ samples (\mthreeneeds\,s at 500\,Hz). Hence the 2\,s hold in decisions D8.
\fix{Version 1 claimed $\sigma_r\approx0.7$\,cm at 0.5\,N\,m and $N = 150$.}
\meta{E.5 [num] (Schur inequality on a Panda Jacobian; numbers from scripts/panda\_m3\_numbers.py)}{(M3)}

\subsection{Angle-only version (stretch)}
\textbf{Claim.} Reading joint deflection under impedance control gives $\sigma_{\text{eff}}^2 = K_p^2\delta^2/12 + \tau_s^2/3$.

\textbf{Derivation.} Quantization is uniform on $[-\delta/2,\delta/2]$, with variance $\delta^2/12$, scaled by $K_p$. Stiction is uniform on $[-\tau_s,\tau_s]$, with variance $\tau_s^2/3$. The two are independent. \hfill $\square$

Both are fixed during a hold and do \emph{not} average down with $N$, so use $2\sigma_{\text{eff}}^2$, not $2\sigma_{\text{eff}}^2/N$, in E.3. The terms are equal at $K_p^\dagger = 2\tau_s/\delta$; below it the variance is within a factor 2 of the friction floor $\tau_s^2/3$. Hand sag is $J_vK_p^{-1}J^{\mathsf T}\bm w$.
\meta{E.6 [sym]}{(stretch)}

% =====================================================================
\section{Belief update and decision}

\subsection{Kalman update of the CoM}
\textbf{Claim.} With prior $\mathcal N(\bm c_0,\Sigma_0)$ and measurement $\bm z = \bm x+\bm v$, $\bm v\sim\mathcal N(0,R)$: $K = \Sigma_0(\Sigma_0+R)^{-1}$, $\bm x^+ = \bm c_0+K(\bm z-\bm c_0)$ and $\Sigma^+ = (I-K)\Sigma_0$.

\textbf{Derivation.}
\begin{enumerate}
  \item Multiply the two Gaussians. Collecting terms in $\bm x$ gives $\Sigma^+ = (\Sigma_0^{-1}+R^{-1})^{-1}$ and $\bm x^+ = \Sigma^+(\Sigma_0^{-1}\bm c_0+R^{-1}\bm z)$.
  \item Rewrite: $\Sigma_0^{-1}+R^{-1} = \Sigma_0^{-1}(\Sigma_0+R)R^{-1}$, so $\Sigma^+ = R(\Sigma_0+R)^{-1}\Sigma_0 = (I-K)\Sigma_0$.
  \item $\Sigma^+R^{-1} = K$ and $\Sigma^+\Sigma_0^{-1} = I-K$. \hfill $\square$
\end{enumerate}
\note{With the real-Panda accuracy (E.5), $R$ is \emph{not} negligible next to $\Sigma_0$ along an object's short axis ($\sigma_0 = 4/\sqrt{12} = 1.15$\,cm for a 4\,cm-wide object). So the prior still matters there.}
\meta{F.1 [num]}{(M3)}

\subsection{Continue or regrasp}
\textbf{Rule.} Continue if $P_{\text{hold}}\ge1-\varepsilon$ ($\varepsilon = 0.1$). Otherwise set the object down, re-rank with the posterior, and regrasp once.

\textbf{Claim.} If $\dstar/\sigma_d$ is large, this is the margin rule $\rho\ge z_{0.9}\sigma_d/\dstar$, with $z_{0.9} = 1.2816$.

\textbf{Derivation.} For $\mu_d\ge0$ the second $\Phi$ term is at most $\Phi(-\dstar/\sigma_d)\approx0$. Then $\Phi((\dstar-\mu_d)/\sigma_d)\ge0.9\iff\mu_d\le\dstar-z\sigma_d$. \hfill $\square$

The rules agree on more than 99.9\% of cases once $\dstar/\sigma_d\ge5$, and on less than 99\% at $\dstar/\sigma_d = 1.5$. The code uses $P_{\text{hold}}$ directly.

\textbf{Through the set-down,} store the CoM in the gripper frame, $\bm c^G = T_G^{-1}\bm c$. Map it back with the release pose, and rotate the covariance the same way: $\Sigma\to R\Sigma R^{\mathsf T}$.
\meta{F.2 [num] (agreement vs $\dstar/\sigma_d$)}{(M3)}

% =====================================================================
\section{Statistics}

\subsection{Confidence interval for a difference in success rates}
\textbf{Claim.} For independent trials, the 95\% half-width is at most $1.96\sqrt{0.5/n}$, which is 12.7 points for $n = 120$. Trials clustered by object ($k$ per object, correlation $\rho$) inflate the variance by $1+(k-1)\rho$ (G.3). With $k = 6$ the half-width is 17.9 points at $\rho = 0.2$ and 23.7 points at $\rho = 0.5$.

\textbf{Derivation.} $\mathrm{Var}(\hat p_1-\hat p_2) = \sum p(1-p)/n\le\frac1{2n}$, since $p(1-p)\le\frac14$. Multiply by the design effect. \hfill $\square$

Intervals are therefore computed by the cluster bootstrap over objects.
\fix{The proposal's ``at worst $\pm13$ points'' ignored clustering.}
\meta{G.1 [arith]; [num] (on simulated clustered data the naive interval undercovers and the cluster bootstrap does not)}{(analysis)}

\subsection{Paired comparison and McNemar's test}
\textbf{Claims.} $\mathrm{Var}\,\bar D = [\pi_b+\pi_c-(\pi_b-\pi_c)^2]/n$. Under $H_0: \pi_b = \pi_c$, $b\mid(b+c)\sim\mathrm{Bin}(b+c,\frac12)$, and $\chi^2 = (|b-c|-1)^2/(b+c)$ approximates it.

\textbf{Derivation.} $ED = \pi_b-\pi_c$ and $ED^2 = \pi_b+\pi_c$. Given that a pair is discordant, it is type $b$ with probability $\pi_b/(\pi_b+\pi_c) = \frac12$. The normal approximation gives $z = (b-c)/\sqrt{b+c}$. \hfill $\square$
\note{McNemar's test needs independent pairs, which clustered trials violate. Use the cluster bootstrap of $\bar D$, or a clustered McNemar test. It also tests ``any difference'', not ``$\ge15$ points''. The rule proposed in decisions D10: a threshold is met when the point estimate reaches it \emph{and} the cluster-bootstrap 95\% interval excludes 0. The paired ``$\pm5$ points'' needs $(\pi_b+\pi_c)[1+9\rho_D]\le0.13$.}
\meta{G.2 [num] (conditional binomial simulated)}{(analysis)}

\subsection{Why resample by object}
\textbf{Claim.} $\mathrm{Var}\,\bar Y = \frac{\sigma^2}{nk}[1+(k-1)\rho]$.

\textbf{Derivation.} Write $Y_{ij} = u_i+e_{ij}$, so $\mathrm{Var}\sum_jY_{ij} = k^2\rho\sigma^2+k(1-\rho)\sigma^2$. Divide by $(nk)^2$ and multiply by $n$. \hfill $\square$
\meta{G.3 [num]}{(analysis)}

\subsection{Fitting the failure point $d_{50}$}
\textbf{Model and gradient.} $P = 1/(1+e^{k(d_{\text{eff}}-d_{50})})$. With $u_i = k(d_{50}-d_i)$ and $\sigma' = \sigma(1-\sigma)$: $\partial\ell/\partial d_{50} = k\sum(y_i-p_i)$ and $\partial\ell/\partial k = \sum(y_i-p_i)(d_{50}-d_i)$.
\note{Simulated outcomes are nearly deterministic in $d_{\text{eff}}$ (R.1). If successes and failures are perfectly separated, the likelihood keeps rising as $k\to\infty$ and no finite maximum exists. So bound $k$, use Firth's penalized likelihood, or report the gap between the largest success and the smallest failure with a bootstrap. On the grid pad $d_{50}$ should be compared with the patch $\dstar$ for the contact that actually occurred, not with a fixed 4.74\,cm.}
\meta{G.4 [num] (gradient; separated data has no finite MLE)}{(analysis)}

% =====================================================================
\section{Hidden-weight objects from scanned meshes}
These are the study's objects (decisions D15, D16). B.2 is the simpler point-mass version used by the calibration rig.

\subsection{The solid and its frame}
\textbf{Definitions.} The solid $\Omega$ is the union of the object's convex collision pieces, sampled on a grid of cubes of side $h$ (2\,mm): a cube belongs to $\Omega$ when its centre lies inside some piece. Piece $j$ is the set $\{\bm p : A_j\bm p + \bm b_j \le \bm 0\}$ given by the outward facet planes of its convex hull. With $N$ cubes at centres $\bm p_i$, the volume is $V = Nh^3$, the centroid is $\bm c = \frac1N\sum_i \bm p_i$, and the long axis $\bm e$ is the first right-singular vector of the centred matrix $[\bm p_i - \bm c]$ (the direction of largest spread). Positions along the axis are $x_i = (\bm p_i - \bm c)\cdot\bm e$, so $\sum_i x_i = 0$.
\meta{H.1 [num] (a box solid's centroid and long axis)}{objects.occupancy, objects.principal\_frame}

\subsection{The heavy-end core}
\textbf{Claim.} Split $\Omega$ into a core $F$ (an end cap beyond a cut plane, minus a 2\,mm skin) and the shell $S = \Omega\setminus F$, each of uniform density, with total mass $M$. With $V_F$, $V_S$ their volumes and $\bar x_F$ the core centroid's position along $\bm e$, the CoM lies at $s$ along $\bm e$ when the core mass is
\begin{equation}\label{eq:core}
m_F(s) = M\,\frac{s\,V_S/\bar x_F + V_F}{V_F + V_S}.
\end{equation}

\textbf{Derivation.}
\begin{enumerate}[nosep]
\item Because $\sum_i x_i = 0$ over the whole solid, $V_F\bar x_F + V_S\bar x_S = 0$, so $\bar x_S = -V_F\bar x_F/V_S$.
\item The CoM along $\bm e$ is $\big(m_F\bar x_F + (M - m_F)\bar x_S\big)/M$. Set it to $s$ and substitute step 1: $Ms = m_F\bar x_F - (M - m_F)V_F\bar x_F/V_S$.
\item Collect $m_F$: $m_F\,\bar x_F(1 + V_F/V_S) = Ms + MV_F\bar x_F/V_S$.
\item Divide by $\bar x_F(1+V_F/V_S) = \bar x_F V/V_S$: $m_F = M(sV_S/\bar x_F + V_F)/V$. \hfill $\square$
\end{enumerate}
\textbf{Checks on the result.} At $s=0$, $m_F = MV_F/V$: both parts have density $M/V$, the uniform object. $m_F$ grows linearly in $s$. Densities: core $\rho_F = m_F/V_F$, shell $\rho_S = (M - m_F)/V_S$.

\textbf{Choosing the cut.} For the largest shift $s_{\max} = 8$\,cm the core must stay physically possible: $\rho_F \le 19\,300$\,kg/m$^3$ (tungsten) and the shell must keep at least 100\,g. Of the cuts on either end that satisfy both, the outermost is used. That gives the smallest, most hidden core, like a hammer head or the batteries at one end of a flashlight. The same $F$ is used at every shift, so the object looks identical; only the two hidden densities change.

\textbf{Sideways CoM.} The core centroid can sit slightly off the axis ($\bar y_F$, $\bar z_F$), which moves the CoM sideways by $(m_F - MV_F/V)\,\bar y_F\,V/(MV_S)$ along $y$ (steps 1--2 with $y$ in place of $x$), and likewise along $z$. The lever arm only depends on the CoM's component perpendicular to the closing axis, and the code computes it from the full 3-D CoM (A.2), so this is carried through rather than assumed zero.
\meta{H.2 [sym] (the algebra), [num] (voxelized CoM equals $s$ at every shift; densities within limits)}{objects.core\_masses, objects.choose\_core}

\subsection{Mass properties from the cubes}
\textbf{Claim.} With cube masses $m_i$ (the core or shell density times $h^3$), CoM $\bm g = \sum m_i\bm p_i/\sum m_i$ and $\bm r_i = \bm p_i - \bm g$, the inertia about the CoM is
\begin{equation}\label{eq:voxinertia}
I = \sum_i m_i\big(|\bm r_i|^2 \mathbb{1} - \bm r_i\bm r_i^{\mathsf T}\big) + \frac{h^2}{6}\Big(\sum_i m_i\Big)\mathbb{1}.
\end{equation}

\textbf{Derivation.}
\begin{enumerate}[nosep]
\item A uniform cube of side $h$ and mass $m$ about its own centre: $I_{xx} = \int (y^2 + z^2)\,dm = m\big(\tfrac{h^2}{12} + \tfrac{h^2}{12}\big) = \tfrac{mh^2}{6}$, the same for every axis, with zero products of inertia by symmetry.
\item Parallel-axis theorem: moving each cube's inertia to $\bm g$ adds $m_i(|\bm r_i|^2\mathbb{1} - \bm r_i\bm r_i^{\mathsf T})$.
\item Sum over cubes. \hfill $\square$
\end{enumerate}
For a shape made of whole cubes the result is exact. For the scanned objects the only error is the 2\,mm sampling of the surface.
\meta{H.3 [num] (exact for a box of whole cubes; within 2\% for a box through the full pipeline)}{objects.mass\_properties}

% =====================================================================
\section{The camera and the detector's frame}

\subsection{From the depth image to points}
\textbf{Claim.} A MuJoCo camera with vertical field of view $\phi$ and an image of $W\times H$ pixels is a pinhole camera with focal length $f = (H/2)/\tan(\phi/2)$ pixels and principal point $(c_x, c_y) = \big(\tfrac{W-1}{2}, \tfrac{H-1}{2}\big)$ in pixel-index coordinates. A pixel $(u, v)$ (column, row from the top-left) with planar depth $z$ is the point
\begin{equation}\label{eq:backproject}
\bm p_{\text{cv}} = \big((u - c_x)\,z/f,\ (v - c_y)\,z/f,\ z\big)
\end{equation}
in the OpenCV camera frame ($x$ right, $y$ down, $z$ forward). The world point is $\bm p_w = \bm o + R_{wc}\,\bm p_{\text{cv}}$ with $R_{wc} = [\,\bm x_{\text{cam}},\ -\bm y_{\text{cam}},\ -\bm z_{\text{cam}}\,]$ (columns), where $\bm o$ and $\bm x_{\text{cam}}, \bm y_{\text{cam}}, \bm z_{\text{cam}}$ are MuJoCo's camera position and axes.

\textbf{Derivation.}
\begin{enumerate}[nosep]
\item MuJoCo's camera looks along $-\bm z_{\text{cam}}$ with $\bm y_{\text{cam}}$ up, and its vertical field of view spans the image height: a point at depth $z$ on the top edge has $y_{\text{cam}}/z = \tan(\phi/2)$ and sits $H/2$ pixels above the centre. So $f = (H/2)/\tan(\phi/2)$, the same horizontally (square pixels).
\item The image spans continuous coordinates $[0, W]\times[0, H]$ with the optical axis at the centre $(W/2, H/2)$. Pixel $(u, v)$ covers $[u, u+1]\times[v, v+1]$, so its centre is at $(u+\frac12, v+\frac12)$, and its offset from the axis is $u + \frac12 - \frac W2 = u - \frac{W-1}{2}$; likewise for $v$.
\item Pinhole projection $u - c_x = f\,x/z$, $v - c_y = f\,y/z$ with $y$ pointing down the image; solve for $x, y$.
\item The OpenCV axes are $\bm x_{\text{cam}}$, $-\bm y_{\text{cam}}$ (down) and $-\bm z_{\text{cam}}$ (forward), which gives the columns of $R_{wc}$. Its determinant is $\det[\bm x, -\bm y, -\bm z] = \det[\bm x, \bm y, \bm z] = 1$, so it is a proper rotation. \hfill $\square$
\end{enumerate}
\textbf{Workspace crop.} The detector gets the region the authors' workspace mask covers: the footprint of everything more than 1\,cm above the calibrated table height, grown by 2\,cm, including the table points inside it. The collision check gets the whole $60\times60$\,cm pick area. No object segmentation from the simulator is used.
\meta{I.1 [mj] (a rendered box: its top face back-projects to the right height and its edges to the right place); [num] ($R_{wc}$ is a rotation)}{scene.camera\_K, scene.backproject, scene.R\_world\_cv}

\subsection{Grasps in the world, and the approach test}
\textbf{Claim.} A detector grasp with rotation $R_c$ (columns: approach, closing direction, finger height) and centre $\bm t_c$ in the camera frame is the world grasp $R_w = R_{wc}R_c$, $\bm t_w = \bm o + R_{wc}\bm t_c$. Its approach is within $\alpha$ of straight down (decisions D3) exactly when $(R_w)_{31} \le -\cos\alpha$.

\textbf{Derivation.} The same frame change as I.1 applied to every column of $R_c$ and to $\bm t_c$. The approach is the first column $\bm a = R_w\bm e_1$; its angle to $-\bm e_z$ is at most $\alpha$ iff $\bm a\cdot(-\bm e_z) \ge \cos\alpha$, i.e.\ $-(R_w)_{31}\ge\cos\alpha$. \hfill $\square$
\meta{I.2 [num] (round trip world $\to$ camera $\to$ world; the test against the angle computed directly)}{scene.grasps\_to\_world, scene.approach\_ok\_world}

\subsection{Is a resting pose stable at every shift?}
\textbf{Claim.} An object at rest on the table stays put only if the ground projection of its CoM lies inside the support polygon, the convex hull of its contact points. With the polygon's edges written as $\bm a_k\cdot\bm q + b_k \le 0$ ($|\bm a_k| = 1$, outward normals), the projection $\bm q$ lies a distance $\delta = -\max_k(\bm a_k\cdot\bm q + b_k)$ inside it (negative if outside).

\textbf{Derivation.} For a convex polygon, $\bm q$ is inside iff it is on the inner side of every edge line. $\bm a_k\cdot\bm q + b_k$ is the signed distance to edge line $k$ (negative inside), so the smallest distance to an edge line is $-\max_k(\cdot)$. Inside a convex polygon, the nearest boundary point lies on the nearest edge line, so this is the distance to the boundary. \hfill $\square$

A resting pose is used only if $\delta\ge5$\,mm at every shift (the hidden core moves the CoM), and the object then has to stay within a few millimetres when simulated for 1\,s at the shifts 0 and 8\,cm.
\meta{I.3 [num] (a square support: inside, on the edge, outside)}{scene.rest\_poses}

\subsection{What a RealSense D435 reports (the depth model)}
The detectors were trained on real D435 depth, so the simulated image gets the D435's main artifacts. The D435 measures depth by active stereo between two infrared imagers a baseline $b = 50$\,mm apart, and reports depth for the left imager in 1\,mm units.

\textbf{Claim 1 (noise grows as $z^2$).} With disparity $\delta = fb/z$ (pixels) measured with error of RMS $\sigma_\delta$ (Intel's typical subpixel figure, 0.08\,px), the depth error has, to first order,
\begin{equation}\label{eq:rsnoise}
\sigma_z = \frac{z^2\,\sigma_\delta}{f\,b}.
\end{equation}
\textbf{Derivation.} $z = fb/\delta$, so $dz/d\delta = -fb/\delta^2 = -z^2/(fb)$; a small error $\epsilon$ in $\delta$ moves $z$ by $-z^2\epsilon/(fb)$, whose RMS is \eqref{eq:rsnoise}. \hfill $\square$
At the pick area ($z\approx0.65$\,m, $f = 927$\,px) this is 0.73\,mm, about the same as the 1\,mm depth unit.

\textbf{Claim 2 (stereo shadows).} The right imager is the left one moved by $b$ along $\bm x_{\text{cam}}$ (rectified pair). A left pixel $(u, v)$ at depth $z$ appears at $(u - fb/z,\ v)$ in the right image. It gets no depth if that position is outside the right image, or if the right image sees a nearer surface there. Consequences: a band of $fb/z$ columns at the left edge is always invalid, as on the real sensor, and behind a step from depth $z_1$ (near) to $z_2$ (far) a strip $fb\,(1/z_1 - 1/z_2)$ pixels wide is invalid.

\textbf{Derivation.} A point $(X, Y, Z)$ in the left OpenCV frame is $(X - b, Y, Z)$ in the right one, so $u_r = f(X-b)/Z + c_x = u - fb/Z$ and $v_r = v$. At a step, the far surface reappears in the right image only past the near edge's position there: the near edge sits at $u_e - fb/z_1$, and a far pixel $u$ maps to $u - fb/z_2$. The far pixels with $u - fb/z_2 < u_e - fb/z_1$ but $u > u_e$ are hidden, a strip of width $fb/z_1 - fb/z_2$. \hfill $\square$

\textbf{Also modelled:} pixels with a depth jump of more than 1.5\,cm inside their $3\times3$ neighbourhood are invalid with probability 0.5 (real stereo matching fails at edges), and the reported depth is rounded to 1\,mm. The noise is seeded per scene, so the same scene always gives the same image.
\meta{I.4 [num] ($\sigma_z$ by Monte Carlo at three depths; left band width), [mj] (the shadow strip behind a rendered box)}{scene.realsense\_depth, scene.stereo\_shadow}

% =====================================================================
\section{Executing a grasp with the arm (B0)}

\subsection{The arm controller}
\textbf{Claim.} The arm's dynamics are $M(q)\ddot q + c(q,\dot q) + g(q) = \tau + \tau_{\text{ext}}$, where $\tau_{\text{ext}}$ collects everything the controller's model leaves out (the held object, contacts). With the control law
\begin{equation}\label{eq:ctrl}
\tau = M(q)\,\ddot q_d + c(q,\dot q) + g(q) + K_p(q_d - q) + K_d(\dot q_d - \dot q),
\end{equation}
the tracking error $e = q_d - q$ obeys $M(q)\,\ddot e + K_d\,\dot e + K_p\,e = -\tau_{\text{ext}}$. So with no external load $e\to0$, and a constant load leaves the steady error $e = -K_p^{-1}\tau_{\text{ext}}$ (the payload ``sag'').

\textbf{Derivation.} Substitute \eqref{eq:ctrl} into the dynamics: $M\ddot q = M\ddot q_d + K_p e + K_d\dot e + \tau_{\text{ext}}$, i.e.\ $M(\ddot q_d - \ddot q) + K_d\dot e + K_p e = -\tau_{\text{ext}}$. With $\tau_{\text{ext}}$ constant and $e$ at rest, $K_p e = -\tau_{\text{ext}}$. Since $M\succ0$ and $K_p, K_d\succ0$ (diagonal, positive), the homogeneous system is exponentially stable (the energy $\tfrac12\dot e^{\mathsf T}M\dot e + \tfrac12 e^{\mathsf T}K_pe$ decreases at rate $\dot e^{\mathsf T}K_d\dot e$ when $M$ varies slowly). \hfill $\square$

\textbf{In the simulator.} $M$, $c$ and $g$ come from the arm-and-hand model; the object is a separate body, so the controller, like the real Panda's, does not know the payload. The gains are libfranka's joint-impedance example ($K_p = \text{diag}(600,600,600,600,250,150,50)$, $K_d = \text{diag}(50,50,50,20,20,20,10)$), and $\tau$ is clipped at the Panda's limits (87\,N\,m on joints 1--4, 12\,N\,m on 5--7). The links have no \texttt{gravcomp} (decisions D12): gravity really acts and the controller really cancels it, so the joint torques carry the payload for M3.
\note{The first version had no $M\ddot q_d$ term. It lagged its own trajectory by 13\,mm at the end of the descent, the fingers closed before the hand arrived and missed the object. With the feedforward term the lag is about 0.1\,mm, and the gripper also waits until the hand is within 1\,mm before closing.}
\meta{J.1 [mj] (holds a pose exactly; the sag under a known payload equals $K_p^{-1}J^{\mathsf T}\bm F$; tracking error during a move)}{robot.Panda.control}

\subsection{From a detector grasp to the hand pose}
\textbf{Claim.} GraspNet's grasp has rotation $R$ with columns approach $\bm a$, closing direction $\bm c$ and finger height, centre $\bm t$ on the object surface, and fingertip depth $\delta$ past $\bm t$. The Panda grasps it with its TCP (the centre between the two pads) at $\bm t + (\delta - h)\,\bm a$, where $h = 8.5$\,mm is half the pad height, and hand axes $[\bm c\times\bm a,\ \bm c,\ \bm a]$. Rotating $R$ by $\pi$ about $\bm a$ ($\bm c\to-\bm c$) gives the same grasp; both are tried.

\textbf{Derivation.} The Panda's fingertip is the far end of its pad, $h$ beyond the pad centre along the approach, so placing the fingertip at depth $\delta$ puts the pad centre at $\delta - h$. The hand's $z$ axis points along the approach and its fingers move along $y$, so $z = \bm a$, $y = \bm c$, and $x = y\times z = \bm c\times\bm a$. This is right-handed: $x\times y = (\bm c\times\bm a)\times\bm c = \bm a(\bm c\cdot\bm c) - \bm c(\bm a\cdot\bm c) = \bm a = z$. The two-finger hand is symmetric under a half turn about its approach axis. \hfill $\square$
\meta{J.2 [num] (orthonormal, $\det = 1$, axes and pad centre as claimed; the half-turn grasp is the same pair of pad positions)}{robot.grasp\_to\_tcp}

\subsection{Velocity and acceleration of the minimum-jerk path}
With $u = t/T$ and $s = 10u^3 - 15u^4 + 6u^5$ (D.1): $\dot s = (30u^2 - 60u^3 + 30u^4)/T$ and $\ddot s = (60u - 180u^2 + 120u^3)/T^2$, by the chain rule ($du/dt = 1/T$). A joint path $q = q_0 + (q_1 - q_0)s$ then has $\dot q = (q_1-q_0)\dot s$ and $\ddot q = (q_1-q_0)\ddot s$, which feed \eqref{eq:ctrl}.
\meta{J.3 [num] (against finite differences)}{robot.min\_jerk\_scale}

% =====================================================================
\appendix
\renewcommand{\thesection}{R}
\section{Results: contact calibration (step 1)}
\textbf{Rig.} A floating Franka hand on a vertical slide, a $20\times4\times4$\,cm box of 0.6\,kg with the hidden weight at $x_h = 9$\,cm, contact settings as in decisions D5--D6. Each trial grasps at the box centre, lifts 10\,cm in 1.2\,s (peak acceleration \calibapeak\,m/s$^2$), and holds 1.5\,s. Rotation is the \emph{tilt} defined in D.3. ``Slip onset'' means a tilt above 1$^\circ$; ``failure'' means above 10$^\circ$ (decisions D2). The point pad uses rigid contacts, so it isolates the friction model. The grid pad is rubber (nominal) or rigid (sensitivity case). Scripts: \texttt{scripts/calibrate\_contact.py} and \texttt{scripts/characterize\_contact.py}.

\subsection{Gripper setups}
\begin{table}[H]
\centering
\caption{Sweep of the hidden shift, 0--8.5\,cm in 0.25\,cm steps (holding-regime stats taken before the first failure)}
\scriptsize
\begin{tabularx}{\textwidth}{@{}>{\raggedright\arraybackslash}X c c c c c c c c@{}}
\toprule
Setup & pts/pad & $f_n$ [N] & depth [mm] & $d{=}0$ & rot.\ before fail [$^\circ$] & holds to [cm] & fails from [cm] & palm catches \\
\midrule
\calibsetups
\bottomrule
\end{tabularx}
\end{table}
Menagerie's gripper drops the box outright. With MuJoCo's default contact settings both pads fail near 4.6\,cm, because creep hides the pad physics. With the calibrated settings, the point pad fails at its soft-finger $\dstar$ and both grid pads near their patch $\dstar$ (\calibgridpred\,cm).

\subsection{Does the onset follow the prediction when the parameters change?}
\begin{table}[H]
\centering
\caption{Measured onsets vs the predicted band [onset at peak lift acceleration, static onset]; deviation from the static prediction in brackets; $\dagger$ = outside the band (bisection resolution 0.004\,cm)}
\scriptsize
\begin{tabular}{@{}lccccccc@{}}
\toprule
pad & $\mu$ & $f_n$ [N] & $a$ [mm] & $m$ [kg] & band [cm] & slip onset, 1$^\circ$ [cm] & failure, 10$^\circ$ [cm] \\
\midrule
\calibgeneral
\bottomrule
\end{tabular}
\end{table}
\textbf{Point pad (rigid contacts).} \calibinbandonePoint\ slip onsets fall in the band. The simulator reproduces the soft-finger model across $\mu$, $f_n$, $a$ and $m$, as it should by construction (A.7).

\textbf{Grid pad, rigid.} \calibinbandoneRigid\ slip onsets fall in the band; the other is 0.7\% above it. The twisting resistance emerges from the spread of the 16 contacts and matches the patch limit surface (A.12) with no fitted parameter. Its 10$^\circ$ failure comes 4--11\% later than the onset, because the object re-sticks.

\textbf{Grid pad, rubber (the study's pad).} The gross failure (10$^\circ$) falls in the band in \calibinbandtenRubber\ cases; the other is 0.2\% above it. Small rotations (1$^\circ$) start about 4--5\% earlier, consistently. This is incipient slip: a compliant pad lets the object creep a little before it gives way, as real soft fingertips do. For the rubber pad the prediction therefore marks the \emph{gross} slip, which is what the study's failure criterion measures.

\textbf{In every case} the 10$^\circ$ failure comes after the slip onset, because the object re-sticks as its lever arm shrinks (A.12 note).

\subsection{Non-ideal grasps (grid pad, rubber and rigid)}
\begin{table}[H]
\centering
\caption{Rubber pad (nominal). One change at a time; prediction from the pad points that touch in that geometry; deviation from the static prediction in brackets; $\dagger$ = outside the band}
\scriptsize
\begin{tabularx}{\textwidth}{@{}>{\raggedright\arraybackslash}X c c c c@{}}
\toprule
Case & pts in prediction & band [cm] & slip onset, 1$^\circ$ [cm] & failure, 10$^\circ$ [cm] \\
\midrule
\calibcharacterizerubber
\bottomrule
\end{tabularx}
\end{table}
\begin{table}[H]
\centering
\caption{Rigid pad (hard plastic, sensitivity case). Same cases}
\scriptsize
\begin{tabularx}{\textwidth}{@{}>{\raggedright\arraybackslash}X c c c c@{}}
\toprule
Case & pts in prediction & band [cm] & slip onset, 1$^\circ$ [cm] & failure, 10$^\circ$ [cm] \\
\midrule
\calibcharacterizerigid
\bottomrule
\end{tabularx}
\end{table}
Read the rubber table at the 10$^\circ$ column (gross slip) and the rigid table at the 1$^\circ$ column (R.2).

\textbf{What the model gets right} (both pads):
\begin{itemize}
  \item a grasp moved along the object (the onset moves by exactly that distance);
  \item a grasp moved up the face;
  \item tilts of 15$^\circ$ and 30$^\circ$ (A.3);
  \item fast and slow lifts (A.6);
  \item a halved timestep.
\end{itemize}

\textbf{Partial contacts.}
\begin{itemize}
  \item With a thin box or a cylinder grasped through its axis, only the two middle rows touch. The two-row patch prediction holds within a few percent on the rubber pad. On the rigid pad, the cylinder slips 28\% earlier, because a rigid flat pad meets a curved surface in a single line.
  \item Pinching a cylinder 10\,mm above its axis tilts the contact normal by 10$^\circ$ and squeezes it downward. Both pads slip 25--35\% earlier than the one-row prediction, which assumes normals along the closing axis.
\end{itemize}

\textbf{What the model misses: misalignment.} When the fingertip is not parallel to the face, it touches along one edge. The two pads then press on diagonally opposite edges, which creates a twisting couple about the vertical axis, and holding that couple uses up friction.
\begin{itemize}
  \item \textbf{Box yawed 3$^\circ$:} the rubber pad fails at about 5.2\,cm ($-31\%$); the rigid pad at about 1.6\,cm.
  \item \textbf{Grasp 8\,mm off-centre sideways:} one finger arrives first and turns the box out of square. Rubber fails at 4.0\,cm, rigid at 1.5\,cm.
  \item \textbf{Box yawed 10$^\circ$:} both fail almost at once. On the rubber pad the box slides toward square, and that sliding spends the friction.
\end{itemize}
These are real effects of parallel jaws. GSNet's in-plane grasp angles come in 15$^\circ$ steps, so misalignments of a few degrees will be common in the study.
\note{The patch model assumes the fingertip is parallel to the object face and the contact normals lie along the closing axis. The full pipeline should estimate, from the point cloud, which pad points will touch and how the surface is oriented under each pad. It can then either correct the grasp's angle or score misaligned grasps with their reduced patch. This is the main modelling task for M2 (decisions D4).}

\end{document}
